% Pillar-5 (MIN-REPORT-RL) results opener
\section{Reading Guide: Evidence Base and Its Provenance}
\label{sec:p5-results}
This reporting proposal is motivated by three internal program observations,
not a causal decomposition of stack and algorithm importance. Exhibit 1 is
an unmatched checkpoint/harness comparison; Exhibit 2 compares label semantics
across different experiments; Exhibit 3 contains twelve configured runs on
one managed stack, three seeds per label. We call these internal program
records Stratum B; the open trainer's toy arithmetic evidence (Stratum C) is
directional only, not an independent external replication.

The earlier pooled benchmark tests, 98-cell inference-grid analyses,
cross-corpus eta-squared ratios, telemetry-item replacement recommendations,
and entropy-based field deprecations are withdrawn from this manuscript.
The current proposal retains all eight reporting requirements, including
base-checkpoint identity. It does not deprecate model or temperature because
a binary-factor entropy is below one bit. A constant label in a sampling-only
corpus provides no test of algorithm importance. No operational CI gate for
retracted Items 13--17 is endorsed.

The argument then proceeds in four steps. \secref{sec:p5-evidence} presents the
exhibits: the label does not pin down the experiment; these observations do
not estimate the relative causal importance of stack and algorithm.
\secref{sec:p5-stack} converts the exhibits
into a standard---the eight-item minimum reportable stack---where each item is
included if and only if we can point to a documented lever through which that
item moved or flipped a result. \secref{sec:p5-threat} generalizes from
exhibits to a threat model: the confound vectors an adversarial (or merely
careless) comparison can ride, graded by flip risk. \secref{sec:p5-toolchain}
shows the standard is cheap to enforce mechanically, which is what
distinguishes a checklist that changes practice from one that decorates
appendices \citep{pineau2020improving}.
